\documentclass[11pt,oneside]{report}
\usepackage[margin=1in]{geometry}
\usepackage[T1]{fontenc}
\usepackage[utf8]{inputenc}
\usepackage{lmodern}
\usepackage{microtype}
\usepackage{amsmath,amssymb}
\usepackage{booktabs,longtable,array}
\usepackage{graphicx}
\usepackage{xcolor}
\usepackage{enumitem}
\usepackage{listings}
\usepackage{fancyhdr}
\usepackage[hidelinks]{hyperref}

\definecolor{frameworkblue}{HTML}{244A73}
\definecolor{frameworkgray}{HTML}{5D6874}
\hypersetup{colorlinks=true,linkcolor=frameworkblue,urlcolor=frameworkblue}
\setlist{nosep}
\setcounter{secnumdepth}{3}
\setcounter{tocdepth}{2}
\setlength{\emergencystretch}{2em}
\pagestyle{fancy}
\setlength{\headheight}{14pt}
\fancyhf{}
\fancyhead[L]{Measured Track-Based Drivetrain Design Framework}
\fancyhead[R]{Methods}
\fancyfoot[C]{\thepage}
\lstset{basicstyle=\ttfamily\small,frame=single,breaklines=true,columns=fullflexible}

\newcommand{\framework}{Measured Track-Based Drivetrain Design Framework}
\newcommand{\bundle}{Track Evidence Bundle}
\newcommand{\code}[1]{\texttt{#1}}
\newcommand{\quality}{\mathcal{Q}}

\title{\textbf{Measured Track-Based Drivetrain Design Framework}\\
       \Large Methods, contracts, uncertainty, and decision reporting\\[1em]
       \large Working Phase 8 checkpoint}
\author{Framework technical documentation}
\date{Version 0.8.0.dev0}

\begin{document}
\pagenumbering{roman}
\maketitle

\begin{abstract}
This document specifies a reproducible method for converting repeated closed-course
GPX recordings into drivetrain design evidence.  It explains the complete chain:
configuration provenance, GPX ingestion, lap reconstruction, centreline and
map-matching, physical feature and response-group contracts, speed-gate confidence,
the versioned \bundle, vehicle and tire equations, obstacle work, bounded and
infinite ideal-CVT mechanisms, gate enforcement, energy closure, uncertainty
representation, paired studies, ranking, attribution, hierarchical reporting, and
provenance.  The document is organized so that the high-level study descriptions in
Chapter~\ref{ch:studies} reference one common set of mechanisms rather than repeating
derivations.  This is important: a study changes which declared inputs are varied; it
does not silently change the physical model.

The current ideal-CVT comparison is one implementation inside a broader measured
track-based framework.  Section~\ref{sec:extensions} states the interfaces by which a
transient CINDER drivetrain, a richer tire law, calibrated obstacle models, or another
driveline may be introduced without rewriting the upstream evidence pipeline.
\end{abstract}

\tableofcontents
\listoftables
\clearpage
\pagenumbering{arabic}

\chapter{Purpose and scope}\label{ch:purpose}

\section{Engineering decision}

The framework addresses a decision that is narrower than full lap-time prediction and
broader than a single CVT tuning calculation:

\begin{quote}
Given reviewed evidence from a measured track, declared vehicle and model uncertainty,
and an explicit comparison contract, which tested drivetrain design is preferable, why,
and how stable is that conclusion?
\end{quote}

The answer has four layers.  First, the track evidence must be credible enough to
constrain the simulation.  Second, the mechanism must close its physical balances and
respect measured gates.  Third, designs must be compared under paired physical
scenarios.  Fourth, the report must distinguish numerical validity, physical output
variation, finite-sample estimation uncertainty, and directional robustness.

The software does not claim that GPS alone identifies suspension work, tire parameters,
driver strategy, or rubber-belt shift mechanics.  Unknowns are represented as explicit
uncertainty or declared limitations rather than hidden fitting freedom.

\section{Current flagship comparison}

The implemented reference problem compares:

\begin{enumerate}
  \item a bounded ideal CVT with declared minimum and maximum reduction ratios; and
  \item an otherwise identical infinite-ratio reference that removes the operating
        ratio window but retains finite launch torque.
\end{enumerate}

The infinite case is a counterfactual opportunity bound.  It is not a realizable
drivetrain and is never reported as a hardware recommendation.  Its purpose is to
isolate what the finite ratio window prevents under the same engine, efficiency,
vehicle, tire, track, gates, driver envelope, and launch-capacity contract.  The exact
powertrain definition appears in Section~\ref{sec:cvt}.

\section{Evidence and model boundaries}

The framework keeps three statements separate:

\begin{description}
  \item[Measured evidence] what was observed in GPX passes, feature surveys, video, or
  another declared source;
  \item[Model declaration] the equations and parameter distributions used to interpret
  the evidence; and
  \item[Decision synthesis] the ranking and warnings produced from completed paired
  cases.
\end{description}

The \bundle{} is the boundary between the first two.  It contains everything a
downstream simulator is allowed to know about the track, but it does not contain the
vehicle model or the final decision.

\chapter{Architecture and reproducible configuration}\label{ch:architecture}

\section{Pipeline}

\begin{center}
\fbox{\begin{minipage}{0.92\textwidth}
\centering
Project and profiles $\rightarrow$ canonical GPX $\rightarrow$ reconstructed laps and
centreline $\rightarrow$ reviewed features and gates $\rightarrow$ \bundle\\[0.4em]
$\rightarrow$ drivetrain/tire/obstacle mechanism $\rightarrow$ paired study execution
$\rightarrow$ physical and uncertainty attribution $\rightarrow$ decision-first report
\end{minipage}}
\end{center}

Each arrow is a contract and a provenance checkpoint.  The mechanism may consume the
bundle, but it may not reinterpret raw GPX internally.  Reporting consumes completed
machine artifacts and cannot change simulation state.

\section{Project ownership}

A project owns track declarations, run identities, vehicle files, studies, and results.
Reusable profiles may provide broad priors or team defaults.  Resolution applies layers
in a declared order and exports the final merged input.  A local override replaces a
leaf but does not erase its source history from the resolution manifest.

Every physical scalar is represented as

\begin{equation}
q = \left(q_0,\; u,\; s,\; \mathcal{D},\; r\right),
\end{equation}

where $q_0$ is the nominal value, $u$ is the unit, $s$ is source provenance,
$\mathcal{D}$ is an uncertainty declaration, and $r$ is the semantic uncertainty role.
A value treated as exact uses a fixed distribution with an explicit reason.  This
prevents a numerical literal from silently masquerading as perfect knowledge.

\section{Units and physical support}

Values are converted to SI during resolution.  Units remain visible in the input
contract and human outputs.  Distribution support is validated before sampling.  The
system rejects, for example, nonpositive mass, efficiency above one, a CVT minimum
reduction greater than its maximum, negative uncertainty scales, invalid probability
vectors, or physically unsupported obstacle parameters.  Rejection is preferred to
silent clipping because clipping changes the declared distribution.

\section{Fingerprints}

Canonical strict JSON serialization produces a content fingerprint

\begin{equation}
h = \operatorname{SHA256}\!\left(\operatorname{canonical}(\mathcal{C})\right),
\end{equation}

where $\mathcal{C}$ is the resolved scientific contract.  Timestamps, process IDs,
machine-dependent paths, and cache counters are excluded from scientific fingerprints.
Study fingerprints own resume workspaces; simulation fingerprints own cached summaries.

\chapter{GPX ingestion}\label{ch:gpx}

\section{Run identity}

Each GPX file has a unique run identifier plus vehicle and driver identifiers.  The
manifest independently controls whether the run contributes to centreline geometry and
gate evidence.  These identities later define a complete paired lap draw in
Section~\ref{sec:paired-gates}; therefore, distinct sessions must not share one run ID.

\section{Safe parsing and canonical telemetry}

The parser uses a hardened XML implementation.  It preserves segment boundaries,
timestamps, coordinates, optional elevation, and optional reported speed.  Invalid
points are diagnosed rather than interpolated invisibly.  Time must be monotonic within
a usable segment.

If no trustworthy speed field exists, speed is derived from horizontal displacement and
elapsed time.  For consecutive points $i-1$ and $i$,

\begin{equation}
v_i = \frac{d_{\mathrm{geo}}(\mathbf{p}_{i-1},\mathbf{p}_i)}
             {t_i-t_{i-1}},
\label{eq:gpx-speed}
\end{equation}

provided the elapsed time is positive and the interval passes diagnostics.  The output
records whether each speed was reported or derived.  Equation~\eqref{eq:gpx-speed} is an
observation channel, not a vehicle dynamic equation.

\section{Diagnostics}

Ingestion reports missing timestamps, coordinate errors, duplicated points, abnormal
time gaps, implausible speed, and derived-speed spikes.  Elevation is retained but is not
differentiated into road grade in the current release.  This limitation is repeated in
Chapter~\ref{ch:limitations} because differentiating noisy altitude would introduce a
new force model and uncertainty contract.

\chapter{Track reconstruction}\label{ch:track}

\section{Local coordinate frame}

Latitude and longitude are projected into a local metric frame $(x,y)$ about a robust
course origin.  The local approximation is appropriate for a compact endurance course;
the original geographic coordinates remain in provenance and review outputs.

\section{Lap detection}

A declared lap-gate event defines $s=0$.  Candidate crossings within a configured
radius are ordered in time.  Adjacent crossings form candidate laps.  A lap is retained
only if timing, moving-speed coverage, maximum-speed, data-gap, and later map-matching
checks pass.  Rejection reasons remain row-level evidence.

This is deliberately not an optimization that selects only the fastest or cleanest
laps.  It is a quality filter whose thresholds are declared before results are viewed.

\section{Centreline construction}

Valid geometry-contributing laps are aligned to the lap gate, resampled, and combined
into a common closed centreline.  Let $\mathbf{c}(s)=[x(s),y(s)]^T$ with
$s\in[0,L)$.  The distance coordinate satisfies

\begin{equation}
s_{j+1}=s_j+\lVert \mathbf{c}_{j+1}-\mathbf{c}_j\rVert_2,
\qquad L=s_{N}.
\end{equation}

The output spacing is a configuration choice, not added physical resolution.  Median
aggregation limits the influence of one poorly aligned pass.  A closed-course seam is
handled explicitly so features and windows may wrap across start/finish.

\section{Map matching}

Each retained GPX point is projected to the centreline, producing along-track coordinate
$s$, signed/unsigned cross-track error, and local quality information.  A lap whose
coverage or map error violates the declared threshold is excluded from behavioral
evidence.  The matched points, not the raw timestamp order alone, define the spatial
windows used for gate metrics.

Centreline curvature is calculated for diagnostics.  The current longitudinal vehicle
model does not spend tire capacity on curvature-derived lateral demand; measured gates
represent repeatable corner-entry behavior.  This boundary prevents an unvalidated
lateral model from being hidden inside track reconstruction.

\chapter{Physical features and response groups}\label{ch:features}

\section{Geometry contract}

A feature declares a stable ID, geographic anchor, horizontal uncertainty, interval
extent, extent uncertainty, source, analysis role, and optional gate candidacy.  After
projection, the bundle records start, anchor, end, wrap status, effective position error,
and review flags.

The anchor is not automatically the physical entry.  The before/after extents express
where the declared physical mechanism is active.  Entry, approach, exit, and recovery
windows are constructed relative to those boundaries.

\section{Response groups}

A response group represents one behavioral speed response.  Multiple nearby physical
features may belong to one response group if drivers treat them as one compound event.
Conversely, separate responses must not be merged only because features are close.  One
accepted speed gate is attached to a response group, while obstacle energy remains
attributable to individual physical features.

This separation prevents double enforcement: a log row and a rut can contribute
separate physical work while sharing one measured entry-speed ceiling.

\section{Pass metrics}

For each valid lap and response group, the framework extracts approach speed, entry
speed, minimum event speed, braking drop, recovery distance, lap median pace, vehicle,
driver, and coordinate quality.  The pass table is the auditable basis of the confidence
score in Chapter~\ref{ch:gates} and paired gate sampling in
Section~\ref{sec:paired-gates}.

\chapter{Speed-gate confidence and enforcement}\label{ch:gates}

\section{Confidence components}

The gate score is an evidence summary, not a posterior probability.  For $n$ eligible
passes and target count $n_{\star}$,

\begin{align}
c_{\mathrm{count}} &= \min\left(1,\frac{n}{n_{\star}}\right),\\
c_{\mathrm{repeat}} &= \max\left(0,1-\frac{\operatorname{IQR}(v_{\mathrm{entry}})}
                                            {s_{\mathrm{repeat}}}\right),\\
c_{\mathrm{brake}} &= \frac{1}{n}\sum_{i=1}^{n}
   \mathbf{1}\!\left(\Delta v_{\mathrm{brake},i}\geq \Delta v_{\star}\right),\\
c_{\mathrm{pace}} &= 1-\left|\operatorname{corr}
   (v_{\mathrm{entry}},v_{\mathrm{lap,median}})\right|,\\
c_{\mathrm{coord}} &= \max\left(0,1-\frac{e_{\mathrm{effective}}}{e_{\max}}\right).
\end{align}

With at least two vehicles, cross-vehicle agreement is

\begin{equation}
c_{\mathrm{vehicle}}=\max\left(0,1-
\frac{\max_j \widetilde v_j-\min_j \widetilde v_j}{s_{\mathrm{vehicle}}}\right),
\end{equation}

where $\widetilde v_j$ is the vehicle-specific median.  A single vehicle receives a
declared neutral score of $0.5$, not fabricated agreement.

The overall score is

\begin{equation}
C=100\sum_k w_k c_k, \qquad \sum_k w_k=1.
\label{eq:gate-score}
\end{equation}

If weights do not sum to one within tolerance, they are normalized and a warning is
emitted.  The report always exports the components, weights, reasons, and recommendation.

\section{Recommendations}

A candidate with coordinate error above the declared limit is \emph{must fix} regardless
of Equation~\eqref{eq:gate-score}.  Too few valid passes requires review.  Otherwise,
declared accept and review thresholds classify the gate.  The score never overrides a
hard geometric failure.

\section{Braking envelope}\label{sec:gate-envelope}

An accepted gate at $s_g$ with target $v_g$ is backward propagated using the effective
braking deceleration $a_b>0$.  At current coordinate $s$, the circular distance ahead is
$d_g(s)$ and the safe speed is

\begin{equation}
v_{\mathrm{safe},g}(s)=\sqrt{v_g^2+2a_b d_g(s)}.
\label{eq:brake-envelope}
\end{equation}

The driver target is the minimum over all active gates.  Effective deceleration is the
minimum of the driver acceleration declaration, brake-force capacity divided by mass,
and straight-line tire brake capacity divided by mass.  A small declared trigger margin
starts the full-brake branch before the mathematical boundary to avoid chatter and
discrete-step overshoot.

Gate enforcement is one-way: it limits excessive simulated entry speed but never
accelerates a weak drivetrain to match a measurement.  Compliance is checked at the
exact gate crossing and permits only the declared numerical tolerance.

\chapter{Track Evidence Bundle}\label{ch:bundle}

\section{Purpose}

The bundle decouples measured-track reconstruction from drivetrain simulation.  It
contains the common centreline, reviewed physical features, response groups, accepted
gate distributions, paired lap identities, obstacle contracts, capability flags,
provenance, schema version, and semantic fingerprint.

It does not contain hidden access to the project, arbitrary Python objects, or an
optimized driver.  A simulator may therefore replay a reviewed evidence snapshot even
when the raw reconstruction pipeline is not rerun.

\section{Version and checksum}

The bundle has a schema version independent of the package version.  A consumer accepts
only a documented compatible family.  Canonical semantic content produces a stable
fingerprint; the serialized file has an adjacent SHA-256 checksum.  The semantic
fingerprint avoids machine-dependent ordering, while the file checksum detects byte
changes.

\section{Capabilities}

Capability flags state whether grade force, lateral coupling, gate sampling identities,
or other optional evidence exists.  Absence is explicit.  In the current release,
elevation may be stored while grade force remains disabled.

\chapter{Obstacle mechanisms}\label{ch:obstacles}

\section{Spatial work contract}

Obstacle models return a local resistance force and optional geometric/traction effects.
Energy is accumulated by the line integral

\begin{equation}
E_{\mathrm{obs}}=\int_{s_0}^{s_1}F_{\mathrm{obs}}(s)\,ds.
\end{equation}

The normalized raised-cosine density over an interval of length $L$ is

\begin{equation}
\phi(x;L)=\frac{1-\cos(2\pi x/L)}{L},\qquad
\int_0^L\phi(x;L)\,dx=1.
\label{eq:raised-cosine}
\end{equation}

It distributes a declared event energy smoothly without changing its integral.

\section{Implemented models}

For a fixed specific event energy $e_f$,

\begin{equation}
F(x)=m e_f\,\phi(x;L).
\end{equation}

For a speed-quadratic event with coefficient $k_v$ and recorded entry speed $v_e$,

\begin{equation}
E=m e_f+k_v v_e^2,\qquad F(x)=E\,\phi(x;L).
\label{eq:speed-obstacle}
\end{equation}

For distributed roughness with specific energy per distance $e_d$,

\begin{equation}
F=m e_d.
\end{equation}

The smooth-profile option uses

\begin{align}
z(x)&=\frac{h}{2}\left[1-\cos\left(\frac{2\pi x}{L}\right)\right],\\
z'(x)&=\frac{h\pi}{L}\sin\left(\frac{2\pi x}{L}\right),\\
z''(x)&=\frac{2\pi^2h}{L^2}\cos\left(\frac{2\pi x}{L}\right),\\
\lambda_N&=\operatorname{clip}\left(1+\frac{v^2z''}{g},
\lambda_{\min},\lambda_{\max}\right).
\end{align}

It can alter local slope, normal-load scale, traction multiplier, and unresolved
dissipative work.  These are reduced-order approximations, not suspension simulation.

\section{Model-form uncertainty}

Alternative obstacle mechanisms are categorical structural uncertainty.  A broad
coefficient in the bundle is not measured lap variation by default.  If repeated event
tests support real event-to-event variation, that quantity may be explicitly assigned
the measured-track role.  The distinction controls which study samples it in
Chapter~\ref{ch:studies}.

\chapter{Vehicle and tire dynamics}\label{ch:vehicle}

\section{State}

The longitudinal state is distance $s$, vehicle speed $v$, and driven-wheel angular
speed $\omega_w$:

\begin{equation}
\dot{s}=v.
\end{equation}

Negative vehicle and wheel speeds are excluded by the integration boundary.

\section{Normal load and tire force}

For modeled grade $\theta$, local normal scale $\lambda_N$, and driven-load fraction
$\gamma$, total and driven normal loads are

\begin{equation}
N=mg\max(\cos\theta,0)\lambda_N,\qquad N_d=\gamma N.
\end{equation}

With surface friction $\mu$, local friction multiplier included in $\mu$, and tire
peak scale $\lambda_\mu$, the force ceiling is

\begin{equation}
F_{\max}=\mu\lambda_\mu N_d.
\end{equation}

Slip speed is

\begin{equation}
v_s=r_w\omega_w-v,
\end{equation}

and the compact saturating tire law is

\begin{equation}
F_t=F_{\max}\tanh\left(\frac{K_s v_s}{F_{\max}}\right).
\label{eq:tire}
\end{equation}

The corresponding nonnegative tire-slip dissipation is

\begin{equation}
P_{\mathrm{slip}}=\max(0,F_t v_s).
\end{equation}

Equation~\eqref{eq:tire} is deliberately compact.  It captures force buildup and a
finite ceiling, but not a complete terrain-dependent tire map.

\section{Resistance and equations of motion}

The implemented forces are

\begin{align}
F_g &= mg\sin\theta,\\
F_{rr} &= C_{rr}N\tanh(v/v_0),\\
F_a &= \tfrac{1}{2}\rho C_DA\,v^2,\\
F_o &= \sum_j F_{o,j}.
\end{align}

The small $v_0$ regularizes rolling-force direction near rest.  Vehicle and wheel
dynamics are

\begin{align}
m\dot v &= F_t-F_g-F_{rr}-F_a-F_o,\label{eq:vehicle}\\
I_w\dot\omega_w &= T_w-T_b-r_wF_t.\label{eq:wheel}
\end{align}

The model also reports curvature-derived lateral demand $m v^2|\kappa|$ as a diagnostic.
It does not reduce $F_{\max}$ because lateral coupling has not yet been validated.

\chapter{Powertrain comparison}\label{ch:powertrain}

\section{Engine}

The current engine is a versioned tabulated torque curve scaled by a declared power
factor.  At engine speed $\omega_e$ and throttle $u\in[0,1]$,

\begin{equation}
T_e=uT_{\mathrm{curve}}(\omega_e),\qquad P_e=T_e\omega_e.
\end{equation}

The ideal CVT targets a declared engine speed $\omega_\star$ associated with target
power $P_\star$.

\section{Bounded ideal CVT}\label{sec:cvt}

Let final-drive reduction be $i_f$ and CVT reduction be $i_c$.  The ratio required to
hold target speed is

\begin{equation}
i_{\mathrm{req}}=\frac{\omega_\star}{i_f\omega_w}.
\end{equation}

The selected bounded ratio is

\begin{equation}
i_c=\operatorname{clip}(i_{\mathrm{req}},i_{\min},i_{\max}).
\label{eq:bounded-ratio}
\end{equation}

Inside the window, the engine remains at target speed.  Below the minimum-reduction
requirement, engine speed rises synchronously with $i_{\min}i_f\omega_w$.  Above the
maximum-reduction requirement, the declared ideal launch clutch may hold the engine at
target while the secondary is slower.

Wheel torque and transmitted power are

\begin{equation}
T_w=\eta i_c i_f T_e,\qquad P_w=\max(0,T_w\omega_w).
\end{equation}

Drivetrain-efficiency loss is $P_e(1-\eta)$.  In launch-clutch mode,

\begin{equation}
P_{\mathrm{clutch}}=\max(0,\eta P_e-P_w).
\end{equation}

Operating shortfall is the counterfactual quantity

\begin{equation}
P_{\mathrm{short}}=\max(0,uP_\star-P_e).
\end{equation}

\section{Infinite-ratio reference}

The reference holds $\omega_e=\omega_\star$ outside a finite operating-ratio window,
but it may not produce infinite launch torque.  It shares the bounded design's launch
cap

\begin{equation}
T_{w,\max}^{\mathrm{launch}}=\eta i_{\max}i_fT_e.
\end{equation}

For nonzero wheel speed, target-power torque is $\eta P_e/\omega_w$.  Therefore

\begin{equation}
T_w^{\infty}=\min\left(\frac{\eta P_e}{\omega_w},
T_{w,\max}^{\mathrm{launch}}\right),
\label{eq:infinite-torque}
\end{equation}

with the launch cap selected directly at zero speed.  Equations
\eqref{eq:bounded-ratio} and \eqref{eq:infinite-torque} define the comparison.  A design
variable may share a cached infinite reference only if it is mathematically absent from
this reference mechanism.  That proof currently permits a limited minimum-ratio reuse;
final drive and maximum ratio are not assumed absent from launch behavior.

\chapter{Integration, events, and energy}\label{ch:numerics}

\section{Time integration}

The solver advances Equations~\eqref{eq:vehicle} and \eqref{eq:wheel} with a declared
fixed integration step and emits a separate report-step trace.  Track completion,
maximum time, invalid state, and other termination causes are explicit.  Feature-entry
and gate-crossing states are interpolated at their exact spatial crossings so a report
sampling interval cannot conceal compliance error.

Short coarse-step runs are useful for software-path checks, but they remain invalid if
the quality tolerances fail.  A production numerical step should be justified by
refinement tests on headline metrics, gate excess, and both balances.

\section{Powertrain energy balance}

Integrated engine-side accounting is

\begin{equation}
E_e=E_w+E_{\eta}+E_{\mathrm{clutch}}+R_p,
\label{eq:powertrain-balance}
\end{equation}

where $R_p$ is the numerical residual.  The reported relative residual divides by a
finite engine-energy scale.

\section{Vehicle energy balance}

Let total kinetic energy include translation and driven-wheel rotation.  Vehicle-side
accounting is

\begin{multline}
E_w+K_0=K_f+E_{\mathrm{slip}}+E_{\mathrm{brake}}+E_{rr}+E_a+E_o\\
+W_g+R_v.
\label{eq:vehicle-balance}
\end{multline}

Grade work $W_g$ is signed.  Current GPX elevation does not contribute to it; a smooth
declared obstacle may contribute a conservative local grade term.  Residual $R_v$ is
normalized by transmitted energy, initial kinetic energy, and the magnitude of grade
work.

\section{Opportunity diagnostics}

Within one case, the diagnostic

\begin{equation}
E_{\mathrm{opp,total}}=E_{\mathrm{clutch}}+E_{\mathrm{short}}
\end{equation}

captures launch and engine operating opportunity.  The bounded-versus-infinite
finite-ratio opportunity loss is the nonnegative difference between the two matched
cases.  It is counterfactual and is \emph{not} added to Equations
\eqref{eq:powertrain-balance} or \eqref{eq:vehicle-balance}.

\chapter{Uncertainty representation and sampling}\label{ch:uncertainty}

\section{Semantic roles}

The distribution shape and engineering meaning are orthogonal.  A stochastic input is
assigned one of:

\begin{description}
  \item[measured track] empirical lap/event variability in the observed track;
  \item[structural] uncertain vehicle, environment, parameter, or model-form
        assumption; and
  \item[initial condition] uncertainty in the declared starting state.
\end{description}

This role separation prevents broad obstacle priors from entering measured-track
robustness merely because the parameters are stored in the bundle.

\section{Marginal distributions}

Numeric marginals may be fixed, normal, truncated normal, uniform, triangular, or
empirical.  Categorical model alternatives use declared discrete probabilities.  Let
$U\sim\mathcal{U}(0,1)$ and $F_i^{-1}$ be the validated inverse marginal transform;
then an independent draw is

\begin{equation}
X_i=F_i^{-1}(U_i).
\end{equation}

The nominal value remains the deterministic baseline even if it differs from the
distribution median.

\section{Correlated groups}

For a declared Gaussian-copula group with correlation matrix $R$,

\begin{align}
\mathbf{Z}&\sim\mathcal{N}(\mathbf{0},R),\\
U_i&=\Phi(Z_i),\\
X_i&=F_i^{-1}(U_i).
\end{align}

The matrix must be symmetric, positive semidefinite, have unit diagonal, reference
known stochastic paths with consistent roles, and not overlap another group.  For
nonnormal marginals, $R$ describes latent Gaussian dependence; it is not generally the
final Pearson correlation.

\section{Paired gate sampling}\label{sec:paired-gates}

When possible, one complete observed run/lap/vehicle/driver identity is drawn and all
gate targets use that lap.  This preserves within-lap pace relationships.  If an
otherwise usable identity lacks one gate, only the missing gate may fall back to its
empirical distribution, and the gate ID is recorded.  Independently combining a fast
corner from one lap with a slow obstacle entry from another is therefore not the default.

\section{Paired designs}

For scenario $j$, all designs receive the same draw $\mathbf{x}_j$.  If $y_{dj}$ is an
output for design $d$, a design contrast uses

\begin{equation}
\Delta_{ab,j}=y_{aj}-y_{bj},
\end{equation}

not a difference between unrelated Monte Carlo samples.  Pairing reduces comparison
noise and makes scenario-level regret interpretable.

\chapter{Study types}\label{ch:studies}

This chapter describes what each run changes.  It relies on the common track, vehicle,
powertrain, gate, energy, and sampling mechanisms defined in Chapters
\ref{ch:track}--\ref{ch:uncertainty}; no study silently substitutes a different equation.

\section{Nominal baseline}

The baseline uses every declared nominal and the selected nominal gate statistic.  It
answers whether the bounded and infinite mechanisms execute and what finite ratio range
costs in one case.  It is a mechanism check, not an uncertainty statement.

\section{Measured-track robustness}

Measured-track robustness samples empirical paired gates and only stochastic inputs with
the measured-track role.  Structural and initial-condition inputs remain nominal.  It
answers whether a conclusion survives plausible versions of the measured course.

This study should not be described as ``all uncertainty in the track.''  Geometry
uncertainty, grade, and obstacle model uncertainty are included only when the formal
contract implements and assigns them appropriately.

\section{Structural sensitivity}

One-at-a-time sensitivity evaluates the exact nominal plus declared marginal quantiles
for each selected numeric input.  Categorical inputs evaluate the nominal and declared
alternatives.  Everything else remains nominal.

For numeric response $y(x)$, local slope near the nominal is estimated from the nearest
levels bracketing $x_0$ when possible.  A global slope uses all tested levels.  The study
answers direction and screening magnitude but omits interactions.

\section{Design sweep}

A sweep evaluates explicit values of one design path under common paired scenarios.
Design inputs are excluded from random sampling in the same study.  It reports both
headline metrics, output bands, paired wins, empirical best fractions, regret, threshold
probabilities, and convergence.  A winning boundary point triggers a recommendation to
extend the domain.

\section{Full uncertainty}

Full uncertainty jointly samples every declared stochastic measured-track, structural,
and initial-condition input, plus gates under the paired policy.  It answers what output
distribution follows from all uncertainty the project currently admits.  It cannot
represent omitted model inadequacy; limitations remain necessary.

\chapter{Decision statistics}\label{ch:decision}

\section{Physical output bands}

For output samples $y_1,\ldots,y_n$, reports include mean, standard deviation, median,
and selected quantiles such as $p_{10}$ and $p_{90}$.  These describe variation across
declared physical scenarios.

\section{Bootstrap estimation intervals}

The framework resamples scenarios with replacement and recomputes statistics.  Percentile
bounds on medians and quantiles describe finite-sample estimation uncertainty.  They are
not a second physical uncertainty distribution.  Paired ranking bootstraps resample
whole scenarios so the within-scenario design relationship is retained.

\section{Wins, best fraction, and regret}

For a lower-is-better metric, design $d$ has empirical best indicator

\begin{equation}
B_{dj}=\mathbf{1}\!\left(y_{dj}=\min_k y_{kj}\right),
\end{equation}

with a documented tie policy.  The reported best fraction is the mean of $B_{dj}$ over
tested scenarios and designs; it is not a universal probability that the hardware is
best outside the declared domain.

Scenario regret is

\begin{equation}
R_{dj}=y_{dj}-\min_k y_{kj}\geq 0.
\end{equation}

Paired win fractions compare two designs on the same scenarios.  Threshold fractions
report the frequency with which a declared engineering limit is exceeded and include a
binomial interval.

\section{Decision gates}

A numerical-quality pass requires all cases to complete, reference dominance, accepted
gate compliance, vehicle energy closure, and powertrain energy closure.  A directionally
robust sweep additionally requires the time and opportunity metrics to select the same
winner, acceptable convergence, and paired win-fraction bootstrap lower bounds above
$0.5$.  Numerical validity and directional robustness are separate output fields.

\chapter{Attribution}\label{ch:attribution}

\section{Physical attribution}

Physical attribution reports where energy went using Equations
\eqref{eq:powertrain-balance} and \eqref{eq:vehicle-balance}.  It includes design-level
bands, scenario-level loss shares, feature-level obstacle work, and closure residuals.
Opportunity diagnostics are displayed separately.

\section{Numeric uncertainty screening}

For numeric input $x_i$ and response $y$, the marginal least-squares slope is

\begin{equation}
\widehat\beta_i=\frac{\sum_j(x_{ij}-\bar{x}_i)(y_j-\bar{y})}
                         {\sum_j(x_{ij}-\bar{x}_i)^2}.
\end{equation}

The normalized elasticity at representative values is

\begin{equation}
\mathcal{E}_i=\widehat\beta_i\frac{\widetilde x_i}{\widetilde y}.
\end{equation}

An uncertainty-weighted effect is

\begin{equation}
W_i=|\widehat\beta_i|s_{x_i},
\end{equation}

and relative screening importance is

\begin{equation}
I_i=\frac{W_i^2}{\sum_k W_k^2}.
\label{eq:importance}
\end{equation}

Slope bootstrap bounds, Pearson association, and Spearman rank association accompany
Equation~\eqref{eq:importance}.  Strong sampled-input correlations generate a warning:
marginal slopes are not causal variance partitions.

\section{Categorical screening}

Categorical alternatives report counts and mean output by level, output span, and a
between-level separation ratio analogous to eta squared.  Sparse levels are flagged.
For structural one-at-a-time studies, category values are explicit alternatives rather
than pretending that model names have numeric quantiles.

\section{Sample-count policy}

Attribution is suppressed below eight scenarios, exploratory from eight through
nineteen, and enabled as screening at twenty or more subject to quality and convergence.
This policy prevents a seven-sample smoke test from producing an authoritative-looking
tornado plot.

\chapter{Convergence and validation}\label{ch:validation}

\section{Monte Carlo convergence}

Reports retain scenario count, Monte Carlo standard error where meaningful, and
split-half median stability.  A status below the screening recommendation explicitly
asks for more replicates.  The question ``are seven samples enough?'' is answered by the
diagnostics rather than a fixed claim: seven can exercise code but cannot establish a
stable distributional decision.

\section{Numerical validation}

The implementation is tested with known limiting cases, exact gate crossings, invalid
support, reference dominance, both energy balances, and timestep refinement.  Feature
obstacle work is reconciled with total obstacle energy.  A coarse integration step may
complete all code paths yet remain invalid for decision.

\section{Execution equivalence}

Scenario results are sorted by replicate and design after parallel execution.  The
scientific artifacts must be identical for serial and parallel workers.  A
content-addressed cache may reduce mechanism-run counts, but cached and uncached
scientific outputs must match.  Manifest timestamps, worker count, and cache counts are
expected operational differences.

\section{Failure publication}

A study writes into a sibling \code{.incomplete} workspace.  Each completed scenario is
checkpointed with the owning study fingerprint.  On success, reports and provenance are
written, checkpoints are removed, and the directory is atomically published.  A failed
run cannot masquerade as a completed result.  Resume rejects changed inputs; restart is
explicit.

\chapter{Reporting and provenance}\label{ch:reporting}

\section{High-to-low hierarchy}

Every result is organized as:

\begin{enumerate}
  \item \code{SUMMARY.md}: finding or recommendation, confidence, warnings, and next
        actions;
  \item \code{decision\_trace.md}: explicit reasoning chain;
  \item \code{REPORT.md}: compact technical evidence and interpretation;
  \item plots: visual location and distribution of effects; and
  \item appendix: complete CSV, JSON, JSONL, bundle, resolved inputs, and traces.
\end{enumerate}

The Markdown report is a regenerable projection of machine artifacts.  Report
regeneration does not rerun physics.  Strict JSON disallows nonstandard NaN and infinity
tokens; unavailable values use a status and null representation.

\section{Provenance graph}

The provenance record connects measured GPX, track build, exact bundle hash, resolved
configuration fingerprint, study fingerprint, framework version, command, and report.
The SVG graph is a human index; the JSON is the machine authority.  The bundle and
resolved inputs make the simulation-facing evidence portable even if original project
paths change.

\chapter{Extension interfaces}\label{sec:extensions}

\section{CINDER drivetrain}

A transient CINDER model can implement the drivetrain adapter while consuming the same
runtime track, driver demand, vehicle state, and scenario inputs.  It must expose
wheel-force capability and explicit channels for clutch work, primary/secondary slip,
shift inertia, actuator work, and any stored drivetrain energy.  It also needs its own
reference and caching policy; the ideal-CVT reference assumptions cannot be copied
without proof.

The comparison hierarchy can then become

\begin{center}
infinite ideal CVT $\leftrightarrow$ bounded ideal CVT $\leftrightarrow$ transient CINDER
CVT,
\end{center}

without rebuilding the measured-track evidence pipeline.

\section{Tire and obstacle models}

A tire adapter maps tire state and normal load to longitudinal force and physical loss.
A brush or empirical model must declare its valid domain, state, sign conventions,
energy channel, and uncertainty inputs.  Obstacle extensions similarly declare geometry,
force/work equations, boundary behavior, physical support, and model-form alternatives.

\section{Other drivetrains}

Any driveline may participate if it can consume the framework state/demand boundary and
return deterministic force, energy, and quality channels.  Upstream GPX and bundle code
should not contain conditionals for a particular driveline family.

\chapter{Limitations and future work}\label{ch:limitations}

\section{Current limitations}

\begin{itemize}
  \item GPX elevation is preserved but not converted to road grade force.
  \item The active vehicle model is longitudinal; curvature is diagnostic and lateral/yaw
        dynamics are absent.
  \item The tire law is reduced-order and does not resolve terrain-dependent combined slip.
  \item Obstacle models are explicit approximations and broad defaults are uncalibrated.
  \item The ideal CVT omits rubber-belt shift transients, pulley forces, thermal behavior,
        and detailed slip mechanics.
  \item The driver follows measured gate ceilings and full-throttle/full-brake logic; it
        is not an optimized human control model.
  \item Uncertainty attribution is screening, not a global causal decomposition.
  \item Joint uncertainty cannot represent assumptions that were never declared.
\end{itemize}

\section{Prioritized future work}

Elevation-to-grade processing requires smoothing, lap alignment, vertical-bias handling,
external terrain comparison, and grade confidence.  Obstacle calibration requires known
geometry, video, vehicle/suspension evidence, and controlled before/after speed tests.
CINDER integration requires a versioned transient drivetrain adapter and expanded energy
state.  Optional expensive global sensitivity may later add Sobol or grouped methods,
especially when interactions materially change a drivetrain decision.

\chapter{Audit checklist}\label{ch:audit}

Before treating a report as engineering evidence, confirm:

\begin{enumerate}
  \item run identities, timestamps, and speed provenance are credible;
  \item lap inclusion and map error were reviewed;
  \item physical feature geometry and response groups are correct;
  \item accepted gates have credible component scores and enough passes;
  \item inherited priors are not described as measurements;
  \item uncertainty roles match their engineering meanings;
  \item the study type answers the actual question;
  \item the design domain includes meaningful alternatives and is extended at a boundary
        optimum;
  \item all numerical quality checks pass;
  \item sample count and convergence support the claimed confidence;
  \item energy attribution is not mixed with counterfactual opportunity loss;
  \item correlation and sparse-category attribution warnings are respected; and
  \item bundle, resolved inputs, seed, fingerprints, and package version are archived.
\end{enumerate}

\appendix
\chapter{Study-to-method cross-reference}

\begin{longtable}{p{0.22\textwidth}p{0.31\textwidth}p{0.37\textwidth}}
\toprule
Run & Varies & Reuses without redefinition \\
\midrule
\endhead
Baseline & No stochastic input & Bundle (Chapter~\ref{ch:bundle}), gates
(Chapter~\ref{ch:gates}), vehicle (Chapter~\ref{ch:vehicle}), powertrain
(Chapter~\ref{ch:powertrain}) \\
Track robustness & Paired gates and measured-track-role inputs & All common mechanisms;
structural inputs remain nominal (Chapter~\ref{ch:uncertainty}) \\
Structural sensitivity & One structural input/alternative at a time & Exact nominal base
and common mechanisms (Chapters~\ref{ch:uncertainty} and \ref{ch:attribution}) \\
Design sweep & Declared design values under paired scenarios & Pairing, decision metrics,
and quality gates (Chapters~\ref{ch:uncertainty} and \ref{ch:decision}) \\
Full uncertainty & All declared stochastic roles jointly & Same physical model and paired
gate policy (Chapters~\ref{ch:uncertainty} and \ref{ch:studies}) \\
\bottomrule
\end{longtable}

\chapter{Output-to-claim cross-reference}

\small
\begin{longtable}{>{\raggedright\arraybackslash}p{0.36\textwidth}>{\raggedright\arraybackslash}p{0.54\textwidth}}
\toprule
Artifact & Claim supported \\
\midrule
\endhead
\code{SUMMARY.md} & concise finding, confidence, warnings, and next action \\
\code{decision\_trace.md} & why the headline follows from evidence and gates \\
\code{REPORT.md} & technical interpretation, compact tables, and limitations \\
\code{replicate\_results.csv} & every scenario/design mechanism comparison \\
\code{scenario\_draws.jsonl} & exact paired inputs and gate identities \\
\code{summary.json} & distributions, bootstrap intervals, ranks, regret, thresholds \\
\code{energy\_accounting.json} & physical partitions and nonadditive diagnostics \\
\path{uncertainty_attribution.json} & screening slopes, associations, categories, warnings \\
\code{convergence.json} & whether scenario count stabilizes reported metrics \\
\code{run\_manifest.json} & numerical quality, counts, seed, cache/resume policy \\
\code{provenance.json} & versioned evidence and configuration lineage \\
\code{track\_bundle.json} & exact simulator-facing measured-track contract \\
\bottomrule
\end{longtable}
\normalsize

\chapter{Companion documentation}

The operational \path{README.md} gives the shortest path.  The companion Markdown files
\path{docs/USER_GUIDE.md}, \path{docs/INPUT_CONTRACT.md},
\path{docs/OUTPUT_REFERENCE.md}, \path{docs/DEVELOPER_GUIDE.md}, and
\path{docs/RELEASE_AND_REPRODUCIBILITY.md} provide field-level and workflow references.
Focused contracts under \code{docs/input}, \code{docs/output}, and \code{docs/methods}
remain the schema-adjacent source for implementers.

\end{document}
